书名中的 ``实战''，核心是
``在编程实践中学习''。本书作为昇腾310B芯片的入门指南，将跳出单纯的理论讲解，通过真实的AI推理部署案例，带读者直观理解昇腾310B的硬件架构特性、Atlas工具链使用逻辑与端侧AI项目开发流程
------
从模型适配、量化优化到推理服务部署，每一个知识点都配套可落地的代码示例，让读者在动手编码的过程中，真正掌握昇腾310B的实战应用能力，实现从
``了解芯片'' 到 ``能用芯片落地项目'' 的跨越。

\subsection*{本书定位与目标读者}\label{本书定位与目标读者}

本书面向以下三类读者：

\begin{itemize}
\tightlist
\item
  \textbf{高校学生 /
  科研新人}：希望通过一套系统化路径快速理解边缘AI硬件与部署流程。
\item
  \textbf{嵌入式 / IoT 工程师}：已有一定Linux / C /
  Python基础，希望把AI模型真正跑在边缘端并做性能调优。
\item
  \textbf{AI应用开发者 /
  创客}：已经能使用主流深度学习框架，希望将训练好的模型迁移到昇腾310B进行高效推理与产品化落地。
\end{itemize}

阅读预期：

\begin{itemize}
\tightlist
\item
  零基础读者可依照``快速起步路径''完成第1\textasciitilde3章+精选案例；
\item
  进阶读者可继续深入算子优化、系统整合与复杂多模型协同部署；
\item
  有项目诉求的团队可参考``方法论 +
  附录模板''直接搭建属于自己的边缘AI应用。
\end{itemize}

\subsection*{学习路径速览（建议路线）}\label{学习路径速览建议路线}

\begin{enumerate}
\def\labelenumi{\arabic{enumi}.}
\tightlist
\item
  环境 + 工具链：硬件认知 → 开发环境装配 → CANN工具初试
\item
  基础模型部署：图像分类→ 目标检测 → 语义分割 / OCR / NLP
\item
  性能优化：模型结构裁剪 → 精度-性能权衡（FP16 / INT8）→ 并行与pipeline
\item
  低级能力：自定义算子 → Profiling → ACL / GE 原理 → 内存与数据通路调优
\item
  系统构建：多进程/多模型协同 → 任务调度 → 监控与日志体系
\item
  实战案例：从需求分析 → 方案设计 → 模型适配 → 部署脚本 → 交付验收
\end{enumerate}

\subsection*{全书结构（初版规划）}\label{全书结构初版规划}

本书分为 \textbf{Part I：理论教程} 与 \textbf{Part II：实验教程}
两大部分。

\subsection{Part I：理论教程}\label{part-iux7406ux8bbaux6559ux7a0b}

{ % do not increment counter
\begin{longtable}[]{@{}
  >{\raggedright\arraybackslash}p{(\linewidth - 6\tabcolsep) * \real{0.2500}}
  >{\raggedright\arraybackslash}p{(\linewidth - 6\tabcolsep) * \real{0.2500}}
  >{\raggedright\arraybackslash}p{(\linewidth - 6\tabcolsep) * \real{0.2500}}
  >{\raggedright\arraybackslash}p{(\linewidth - 6\tabcolsep) * \real{0.2500}}@{}}
\toprule\noalign{}
\begin{minipage}[b]{\linewidth}\raggedright
章节
\end{minipage} & \begin{minipage}[b]{\linewidth}\raggedright
标题
\end{minipage} & \begin{minipage}[b]{\linewidth}\raggedright
内容聚焦
\end{minipage} & \begin{minipage}[b]{\linewidth}\raggedright
读者收益
\end{minipage} \\
\midrule\noalign{}
\endhead
\bottomrule\noalign{}
\endlastfoot
Chapter 1 & 昇腾310B边缘计算基础 & 边缘计算概念、竞品对比、310B规格 &
建立边缘计算认知与选型能力 \\
Chapter 2 & CANN 软件栈核心组件解析 & 异构架构、ATC转换工具、msame推理 &
掌握CANN基础与模型转换流程 \\
Chapter 3 & 昇腾 PyTorch 扩展迁移基础 &
环境搭建、torch\_npu插件、经典网络迁移 & 学会在昇腾上运行PyTorch模型 \\
Chapter 4 & PyACL 应用开发基础 & ACL运行时、模型推理流程、DVPP媒体处理 &
掌握C/C++底层应用开发能力 \\
Chapter 5 & 算子开发基础 & TBE DSL流程、310B计算架构、自定义算子 &
理解底层计算原理，具备算子扩展能力 \\
Chapter 6 & 系统工程与高可用部署 &
Profiling分析、AIPP硬件加速、多级流水线 &
掌握性能调优与系统工程化技能 \\
Chapter 7 & 项目实战方法论与交付模板 &
需求分析、评测体系、标准化交付流程 & 建立规范的AI项目交付方法论 \\
Chapter 8 & 综合实战案例集 & 案例规范、目录结构、典型案例概览 &
了解实战案例的组织与核心逻辑 \\
Chapter 9 & 附录与工具箱 & 错误速查、参数模板、Checklist、FAQ &
提供高效查阅的工具与参考资料 \\
\end{longtable}
}

\subsection{Part II：实验教程}\label{part-iiux5b9eux9a8cux6559ux7a0b}

包含从 \textbf{Case 0} 到 \textbf{Case 9}
的十个动手实验，涵盖环境搭建与九大端侧AI实战案例（人脸打卡机、实时跟踪、智能电子琴、掌纹识别、数据采集仪、智能小车、智能相册、手势识别、聊天机器人）。

\subsection*{各模块核心要点概述}\label{各模块核心要点概述}

\textbf{1. 昇腾310B边缘计算基础} - 边缘计算 vs
云计算：概念辨析与应用场景 - 边缘计算卡生态：NVIDIA、华为、瑞芯微对比 -
昇腾310B 技术规格详解

\textbf{2. CANN 软件栈核心组件解析} - CANN 异构计算架构与 CUDA 对比 -
ATC 模型转换工具流程、参数详解与注意事项 - msame 推理工具使用概览

\textbf{3. 昇腾 PyTorch 扩展迁移基础} - PyTorch (torch\_npu)
环境安装与配置 - 基础网络实现：Linear, MLP, LeNet - 现代 CNN
移植实战：AlexNet, VGG, ResNet

\textbf{4. PyACL 应用开发基础} - 快速入门：Hello World 与运行模式 -
初始化与运行时管理：Device, Context, Stream, Memory -
模型推理完整流程：加载、Dataset准备、执行 - DVPP 图像/视频处理能力
(JPEGD, VPC)

\textbf{5. 算子开发基础} - 为什么需要自定义算子：TBE vs AICPU - 310B
硬件架构：存储层级与向量计算单元 - TBE DSL 开发流程与 LogSoftmax 实例

\textbf{6. 系统工程与高可用部署} - 性能优化全景：木桶效应与 Amdahl 定律
- Profiling 性能分析：采集与关键视图解读 -
核心优化技术：AIPP、零拷贝、多级流水线 -
实战演练：车辆检测系统的极致优化之路

\textbf{7. 项目实战方法论与交付模板} - 需求澄清 Canvas 与指标分层 -
Baseline 策略、评测集设计与迭代看板 - 资产沉淀、交付目录与上线 Checklist
- 验收、回归与风险管理

\textbf{8. 综合实战案例集} - 案例统一模板与目录结构规范 -
实战案例概览与重点分析 (人脸、跟踪、音频等) -
结果记录、差异报告与自动化保障

\textbf{9. 附录与工具箱} - 常见报错速查表 - 模型转换参数模板合集
(ResNet, YOLO, INT8) - 性能与质量 Checklist - 术语表、FAQ 与推荐资源

\textbf{10. 实验教程 (Case 0 - 9)} - \textbf{Case
0}：开发板硬件介绍、系统安装 (刷写/启动/网络)、环境验证 - \textbf{Case
1}：智能打卡机 (Web应用, 人脸识别) - \textbf{Case 2}：边缘端实时目标跟踪
(检测+关联) - \textbf{Case 3}：智能电子琴 (手势识别+音频合成) -
\textbf{Case 4}：智能掌纹识别机 (预处理, 特征提取) - \textbf{Case
5}：智能数据采集仪 (传感器, 数据分析) - \textbf{Case 6}：智能小车 (视觉,
路径规划, 运动控制) - \textbf{Case 7}：智能相册 (识别, 分类, 检索) -
\textbf{Case 8}：手势识别 (数据采集, 实时系统) - \textbf{Case
9}：智能聊天机器人 (LLM, 语音交互, RAG)

\subsection*{实践驱动与开源协作}\label{实践驱动与开源协作}

本书所有示例代码、脚本、案例与附录工具均开源托管于本仓库。欢迎通过 Issue
/ PR 反馈问题、提交改进、补充案例或翻译。我们鼓励： - 增补新模型 /
新任务的部署范式 - 分享自定义算子优化经验 - 提交性能测试报告（含硬件信息
+ 指标） - 翻译与文档校对

\subsection*{如何使用本书}\label{如何使用本书}

{ % do not increment counter
\begin{longtable}[]{@{}
  >{\raggedright\arraybackslash}p{(\linewidth - 6\tabcolsep) * \real{0.2500}}
  >{\raggedright\arraybackslash}p{(\linewidth - 6\tabcolsep) * \real{0.2500}}
  >{\raggedright\arraybackslash}p{(\linewidth - 6\tabcolsep) * \real{0.2500}}
  >{\raggedright\arraybackslash}p{(\linewidth - 6\tabcolsep) * \real{0.2500}}@{}}
\toprule\noalign{}
\begin{minipage}[b]{\linewidth}\raggedright
读者类型
\end{minipage} & \begin{minipage}[b]{\linewidth}\raggedright
推荐阅读路径
\end{minipage} & \begin{minipage}[b]{\linewidth}\raggedright
目标
\end{minipage} & \begin{minipage}[b]{\linewidth}\raggedright
补充建议
\end{minipage} \\
\midrule\noalign{}
\endhead
\bottomrule\noalign{}
\endlastfoot
零基础学生 & Chapter 1 → Chapter 2 → Case 0 → Case 1 & 能跑通首个模型 &
结合案例做改动实验 \\
嵌入式工程师 & Chapter 4 → Chapter 5 → Chapter 6 & 掌握底层开发与优化 &
关注资源管理与算子开发 \\
AI应用开发者 & Chapter 2 → Chapter 3 → 选读 Case 案例 & 快速场景落地 &
重点关注 PyTorch 迁移与部署 \\
技术负责人 & Chapter 1 → Chapter 7 → Chapter 8 & 构建团队方法论 &
制定内部交付标准与模板体系 \\
\end{longtable}
}

\subsection*{更新与版本计划}\label{更新与版本计划}

\begin{itemize}
\tightlist
\item
  v0.1（当前）：结构规划 + Chapter 1\textasciitilde3 初稿 + Case
  0\textasciitilde1 示例
\item
  v0.3：补齐核心部署链路 (Chapter 4) + 性能调优初稿 (Chapter 6)
\item
  v0.6：全案例上线 (Case 0\textasciitilde9) + 工程化章节完善
\item
  v1.0：补齐附录 (Chapter 9) + 全面审校 + PDF / LaTeX 发行
\end{itemize}

\subsection*{许可证与引用}\label{许可证与引用}

本书内容采用 Apache 2.0 许可证。引用本书内容请注明： \textgreater{}
《昇腾310B实战：从入门到精通边缘计算与人工智能》（GitHub:
zhouxzh/Ascend310）

\begin{center}\rule{0.5\linewidth}{0.5pt}\end{center}

欢迎加入共建，一起把``边缘AI实战''这件事做成！
